%2multibyte Version: 5.50.0.2960 CodePage: 65001
% MQF_Manuale_Master_SW55_Memoir.tex
% Master del manuale MQF costruito sullo shell SW5.5 "Memoir Class for Configurable Typesetting".
% Versione modulare: i capitoli scritti sono inclusi con \input.%
% POSIZIONE DEI FILE
% Questo master deve stare nella cartella 01_Manuale.
% I capitoli devono stare nella sottocartella 01_Manuale/Capitoli.
% Non mettere il master dentro Capitoli.
% I capitoli futuri restano come traccia, ma non generano pagine nel documento.
% Evita pagine bianche dovute all'apertura dei capitoli su pagina dispari.
% ------------------------------------------------------------
% Nomi italiani
% ------------------------------------------------------------
% ------------------------------------------------------------
% Macro matematiche ufficiali del progetto MQF
% Coerenti con MQF_Notazione.tex
% ------------------------------------------------------------
% ------------------------------------------------------------
% Ambienti di tipo teorema.
% I nomi interni restano compatibili con lo shell SW.
% Le etichette stampate sono in italiano.
% ------------------------------------------------------------

%TCIDATA{OutputFilter=latex2.dll}
%TCIDATA{Version=5.50.0.2960}
%TCIDATA{LaTeXparent=1,1,MQF_Manuale_Master.tex}


% Capitoli/MQF_Cap_06_Processi_stocastici_in_tempo_continuo.tex

\chapter[Processi stocastici, Tempo continuo]{Processi stocastici in tempo continuo: diffusioni, mean-reversion e simulazione}

\label{cap:tempo-continuo}

\section{Obiettivi della lezione}

Una banca deve stabilire oggi il prezzo equo di un contratto il cui
valore dipender\`{a} dall'andamento di un titolo nei prossimi mesi;
un fondo pensione deve valutare un'obbligazione i cui flussi sono
legati all'inflazione futura e ai tassi di interesse. In entrambi i
casi il valore di oggi dipende dall'intera traiettoria futura di una
o pi\`{u} grandezze incerte. Con quali modelli si descrive
l'evoluzione continua di prezzi, tassi e inflazione, e come si
generano le traiettorie su cui fondare queste valutazioni?

Il capitolo precedente ha introdotto il linguaggio dei processi
stocastici --- traiettoria, momenti, filtrazione, martingala, scenario
--- e lo ha applicato a un modello in tempo discreto. Tutti quegli
strumenti, lo abbiamo sottolineato, sono stati definiti in modo da
valere per qualunque insieme di indici temporali. In questo capitolo
li mettiamo all'opera nel contesto del \emph{tempo continuo}, in cui
l'indice $t$ varia con continuit\`{a} su un intervallo e le traiettorie
del processo diventano funzioni del tempo.

Il tempo continuo \`{e} il linguaggio in cui sono formulati i
principali modelli della finanza quantitativa: i prezzi delle
attivit\`{a}, i tassi di interesse, i tassi di inflazione e gli altri
fattori di rischio vengono descritti, nella pratica di mercato, da
\emph{equazioni differenziali stocastiche}. La ragione non \`{e}
soltanto formale. La formulazione continua consente di costruire
modelli con poche grandezze interpretabili --- una tendenza di fondo,
una volatilit\`{a}, una velocit\`{a} di richiamo verso un livello di
equilibrio --- e di simularne le traiettorie su una griglia temporale
fitta, generando gli scenari su cui si fonda la valutazione di prodotti
finanziari complessi.

Questo capitolo presuppone la familiarit\`{a}, acquisita in corsi
precedenti, con il moto browniano, il calcolo stocastico e il lemma di
It\^{o}. Tali strumenti verranno richiamati quando necessario, ma non
sviluppati: l'obiettivo non \`{e} la costruzione della teoria del
calcolo stocastico, bens\`{\i} l'uso dei modelli che su di essa si
fondano per descrivere e simulare l'evoluzione dei fattori di rischio
finanziari.

La trattazione procede dal generale al particolare. Si parte dal
quadro unificante delle equazioni differenziali stocastiche, si
introducono poi i modelli specifici di uso pi\`{u} comune --- il moto
browniano geometrico per i prezzi azionari, i processi mean-reverting
di Ornstein--Uhlenbeck e di Cox--Ingersoll--Ross per i tassi e
l'inflazione --- e si passa infine ai sistemi multivariati, in cui
pi\`{u} fattori di rischio evolvono congiuntamente con una struttura
di correlazione. L'intero percorso \`{e} orientato a fornire la base
teorica per la lezione applicativa in Python, in cui questi modelli
verranno simulati e impiegati nella valutazione di prodotti
finanziari.

Al termine del capitolo lo studente dovr\`{a} essere in grado di:

\begin{itemize}

\item interpretare un'equazione differenziale stocastica nella forma
$dX_t=a(X_t,t)\,dt+b(X_t,t)\,dW_t$, riconoscendo il ruolo del termine
di tendenza (drift) e del termine di diffusione;

\item descrivere il moto browniano geometrico, ricavarne la soluzione
e i momenti mediante il lemma di It\^{o}, e interpretarne i parametri,
riconoscendone al contempo i limiti come modello per tassi e
inflazione;

\item distinguere i processi mean-reverting di Ornstein--Uhlenbeck e
di Cox--Ingersoll--Ross, comprenderne la struttura di richiamo verso
un livello di equilibrio e la differenza cruciale relativa alla
positivit\`{a} delle traiettorie, con la condizione di Feller;

\item costruire un sistema di diffusioni multivariato in cui pi\`{u}
fattori evolvono con una struttura di correlazione assegnata, e
comprendere come tale correlazione si traduca operativamente nella
simulazione;

\item applicare lo schema di discretizzazione di Eulero--Maruyama per
generare traiettorie simulate dei modelli del capitolo, collegando la
simulazione alla stima Monte Carlo introdotta nel Capitolo~5.

\end{itemize}

Il capitolo fornisce la base teorica diretta per la lezione applicativa
in Python (Lezione~7), nella quale i modelli qui presentati verranno
implementati per la simulazione di traiettorie e per la valutazione di
prodotti finanziari il cui valore dipende dall'intera evoluzione dei
fattori di rischio sottostanti.

